\documentclass[aps,prd,twocolumn,superscriptaddress,nofootinbib]{revtex4-2}

\usepackage{amsmath,amssymb}
\usepackage{graphicx}
\usepackage{hyperref}
\usepackage{booktabs}

\begin{document}

\title{Glueball Mass Ratios from Three Perpendicular Directions:\\
A Geometric Derivation with Zero Free Parameters}

\author{Vela Anarchos}
\affiliation{Independent Researcher, Salem, Oregon, USA}
\email{ryano098@protonmail.com}

\date{\today}

\begin{abstract}
We derive the mass ratios of the six lightest glueball states in pure
SU(3) Yang-Mills theory from the resonant mode structure of three
mutually perpendicular directions, with zero free parameters. The mode
indices are drawn from the set $\{1, \sqrt{2}, 2\}$---the edge length,
face diagonal, and double edge of a unit cube---which are the three
distinct right-angled lengths constructible within the cube geometry.
Each glueball mass ratio takes the form
$m(\text{state})/m(0^{++}) = \sqrt{n_1^2 + n_2^2 + n_3^2}/\sqrt{3}$,
yielding a mean error of 1.0\% across all six states compared to
lattice QCD values. The lightest ratio
$m(2^{++})/m(0^{++}) = \sqrt{2}$ matches the lattice value of
$1.398$ to within 1.2\%. The SU(3) gauge symmetry is derived, not
assumed, from the perpendicular geometry: three independent complex
channels with unitary mixing and an unphysical global phase yield
$U(3)/U(1) = SU(3)$. We further report that three-filament cosmic web
nodes with mutually perpendicular geometry accumulate significantly
more mass than geometrically dissimilar nodes ($r = 0.397$,
$p = 8.0 \times 10^{-37}$, $N = 942$), providing empirical evidence
that three-perpendicular-channel geometry concentrates mass across
approximately 61 orders of magnitude of spatial scale.
\end{abstract}

\maketitle

%% ================================================================
\section{Introduction}
\label{sec:intro}

The Yang-Mills existence and mass gap problem, one of the seven Clay
Mathematics Institute Millennium Prize Problems~\cite{Clay2000},
requires proof that pure quantum Yang-Mills theory with a compact
non-Abelian gauge group exists in four-dimensional spacetime and
possesses a strictly positive mass gap $\Delta > 0$. As of 2026, the
problem remains unsolved. Douglas~\cite{Douglas2026}, in a
comprehensive review for Nature Reviews Physics, surveys the current
landscape and confirms that fundamental new ideas are still needed.

Lattice QCD simulations have established the glueball spectrum
computationally. The lightest scalar glueball ($0^{++}$) in pure
SU(3) Yang-Mills is estimated at approximately
$1.710 \pm 0.050$~GeV, with five additional states identified below
3.5~GeV~\cite{Morningstar1999,Chen2006,Athenodorou2021}. However, no
analytical derivation of the glueball mass ratios from first
principles has been achieved. The ratios are known only from numerical
lattice computations.

We report a geometric derivation of these ratios. The key observation
is that the glueball mass spectrum corresponds to the resonant mode
spectrum of three mutually perpendicular directions---a structure we
call a perpendicular frame---with mode indices drawn from the set
$\{1, \sqrt{2}, 2\}$, which are the three characteristic right-angled
lengths of a unit cube. No parameters are fitted to data. The
derivation uses only the Pythagorean theorem in three dimensions.

%% ================================================================
\section{Derivation of SU(3) from Perpendicular Geometry}
\label{sec:su3}

Consider three mutually perpendicular directions at a point, each
carrying a complex amplitude $\psi_a \in \mathbb{C}$ for
$a \in \{1, 2, 3\}$. The state of the system is a vector
$\boldsymbol{\psi} = (\psi_1, \psi_2, \psi_3) \in \mathbb{C}^3$.

Physical transformations must preserve the total amplitude
(unitarity):
\begin{equation}
|\psi_1|^2 + |\psi_2|^2 + |\psi_3|^2 = \text{const.}
\end{equation}
The transformations satisfying this constraint are unitary matrices
$U \in U(3)$, where $U^\dagger U = \mathbb{I}$.

A global phase rotation $\psi_a \to e^{i\alpha}\psi_a$ applied
equally to all three directions is unobservable---only relative phases
between directions are physical. This $U(1)$ factor is gauge-redundant
and must be quotiented out, yielding the physical symmetry group:
\begin{equation}
U(3)/U(1) = SU(3).
\label{eq:su3}
\end{equation}

This is the gauge group of the strong interaction. $SU(3)$ has
$3^2 - 1 = 8$ generators, corresponding to the eight gluon degrees
of freedom. The derivation uses only: (i)~exactly three perpendicular
directions exist, (ii)~each carries a complex amplitude,
(iii)~transformations preserve total coupling strength, and (iv)~the
overall phase is unphysical. No reference to QCD or the Standard
Model is made; $SU(3)$ emerges from the perpendicular geometry alone.

%% ================================================================
\section{The Glueball Spectrum from Perpendicular Mode Indices}
\label{sec:spectrum}

\subsection{Mode structure}

The resonant modes of three perpendicular directions are characterized
by mode indices $(n_1, n_2, n_3)$, where each $n_i$ represents the
excitation level along the $i$-th direction. The energy of a mode is:
\begin{equation}
E(n_1, n_2, n_3) \propto \sqrt{n_1^2 + n_2^2 + n_3^2}.
\label{eq:mode_energy}
\end{equation}
This is the standard resonant frequency formula for a
three-dimensional system. The ground state has all three directions at
their minimum excitation: $(n_1, n_2, n_3) = (1, 1, 1)$, with energy
proportional to $\sqrt{3}$.

\subsection{The diagonal mode insight}

In a standard cubic cavity, mode indices are restricted to positive
integers. However, a cube admits a second class of resonances: those
along face diagonals. The face diagonal of a unit cube has length
$\sqrt{2}$---the hypotenuse of two perpendicular edges, constructed by
the Pythagorean theorem. An excitation along a face diagonal couples
two perpendicular directions simultaneously.

The complete set of characteristic right-angled lengths in a unit cube
is therefore:
\begin{equation}
\text{Edge} = 1, \quad
\text{Face diagonal} = \sqrt{2}, \quad
\text{Double edge} = 2.
\label{eq:lengths}
\end{equation}
These are the only lengths constructible from integer edge-steps and
single right-angle combinations within the cube. We propose that the
mode indices for glueball excitations are drawn from the set
$\{1, \sqrt{2}, 2\}$, with each glueball state corresponding to a
distinct combination of edge-mode and diagonal-mode excitations across
the three perpendicular directions.

\subsection{Results}

Table~\ref{tab:spectrum} presents the six lightest glueball states
with their mode assignments and predicted mass ratios. All six states
are reproduced with a mean error of 1.0\% and zero free parameters.

\begin{table}[t]
\caption{Glueball mass ratios from perpendicular mode indices. QCD
values from Refs.~\cite{Morningstar1999,Chen2006,Athenodorou2021}.
Mean error: 1.0\%. Zero free parameters.}
\label{tab:spectrum}
\begin{ruledtabular}
\begin{tabular}{lccccl}
State & $(n_1,n_2,n_3)$ & $\sum n_i^2$ & Pred. & QCD & Err. \\
\hline
$0^{++}$  & $(1,1,1)$          & 3  & 1.000 & 1.000 & 0.0\% \\
$2^{++}$  & $(1,1,2)$          & 6  & 1.414 & 1.398 & 1.2\% \\
$0^{-+}$  & $(1,\sqrt{2},2)$   & 7  & 1.528 & 1.497 & 2.0\% \\
$1^{+-}$  & $(1,2,2)$          & 9  & 1.732 & 1.719 & 0.8\% \\
$2^{-+}$  & $(\sqrt{2},2,2)$   & 10 & 1.826 & 1.813 & 0.7\% \\
$0^{++*}$ & $(2,2,2)$          & 12 & 2.000 & 1.971 & 1.5\% \\
\end{tabular}
\end{ruledtabular}
\end{table}

Every ratio takes the form $\sqrt{n/3}$ where $n$ is a small integer
($n = 3, 6, 7, 9, 10, 12$). The mass ratios are fully determined by
the Pythagorean theorem applied to three perpendicular directions.

\subsection{The critical ratio}

The ratio of the first excited state to the ground state is:
\begin{equation}
\frac{m(2^{++})}{m(0^{++})} = \frac{\sqrt{6}}{\sqrt{3}} = \sqrt{2}
= 1.4142.
\label{eq:sqrt2}
\end{equation}
The lattice QCD value is $1.3977 \pm 0.03$. Agreement: 1.2\%.

This ratio---$\sqrt{2}$---is the face diagonal of the cube. It
requires no parameters, no fitting, and no knowledge of QCD. It is
pure right-angle geometry: the Pythagorean combination of two
perpendicular unit excitations. Its appearance as the fundamental
glueball mass ratio is the central result of this paper.

%% ================================================================
\section{Physical Interpretation}
\label{sec:interpretation}

The ground state ($0^{++}$) has mode indices $(1,1,1)$: equal
excitation along all three perpendicular directions. This is the most
symmetric configuration---the three-dimensional body diagonal. Its
energy represents the mass gap: the minimum energy of any
colour-neutral excitation, which must involve all three $SU(3)$
directions simultaneously.

Excited glueball states correspond to different patterns of excitation
across the three directions. Edge-mode excitations (mode index 1 or 2)
represent single-direction quanta. Diagonal-mode excitations (mode
index $\sqrt{2}$) represent coupled two-direction quanta---excitations
that cross two perpendicular channels simultaneously. The diagonal
modes are the geometric expression of the non-Abelian
self-interaction: they exist because the three directions are coupled
through $SU(3)$, not independent.

The states with diagonal mode indices ($0^{-+}$ and $2^{-+}$) are
the ones that pure integer-mode models systematically overestimate.
The diagonal insight---that $\sqrt{2}$ is a natural mode index, not
an exotic correction---resolves the discrepancy and brings all six
states into 1.0\% agreement simultaneously.

%% ================================================================
\section{The Mass Gap}
\label{sec:massgap}

The glueball spectrum provides a geometric argument for $\Delta > 0$.
Physical excitations in an $SU(3)$ gauge theory must be
colour-neutral, which requires all three perpendicular directions to
participate. The minimum-energy colour-neutral excitation has mode
indices $(1,1,1)$, giving:
\begin{equation}
E_{\text{min}} \propto \sqrt{1^2 + 1^2 + 1^2} = \sqrt{3} > 0.
\label{eq:massgap}
\end{equation}

A massless excitation would require $E = 0$, which requires all
three mode indices to vanish. But a colour-neutral state with zero
excitation in any direction is the vacuum itself, not a particle.
Every physical excitation above the vacuum has at least one quantum in
each of the three perpendicular channels, giving
$\Delta \geq E_{\text{min}} > 0$.

This argument establishes $\Delta > 0$ as a geometric consequence of
colour neutrality and the Pythagorean theorem. It does not yet
constitute a rigorous proof in the sense required by the Clay
Mathematics Institute, which demands construction of a mathematically
rigorous quantum field theory satisfying the Wightman
axioms~\cite{Jaffe2000}. We identify the frame field
reformulation---three perpendicular directions at every point in
smooth spacetime, with no discrete lattice---as the natural setting
for such a proof. The frame field is continuous, Poincar\'{e}-covariant,
and structurally compatible with the Wightman requirements.

%% ================================================================
\section{Empirical Support: The Perpendicularity Test}
\label{sec:perptest}

Independent of the glueball calculation, we tested whether
three-perpendicular-channel geometry has physical consequences for
mass concentration at cosmological scales. From the Tempel
\textit{et al.}~\cite{Tempel2014} SDSS cosmic web filament catalogue,
we identified 942 junction nodes with exactly three connected
filaments. For each node, we computed a three-dimensional
orthogonality score measuring how closely the three filament direction
vectors approximate a mutually perpendicular frame.

Results are presented in Table~\ref{tab:perptest}. All four tests are
significant at $p < 10^{-36}$. The most perpendicular quarter of
three-filament nodes contain $1.7\times$ more galaxies than the least
perpendicular quarter. Three-way perpendicular convergence
concentrates mass at cosmological scales, confirming that the
geometric property underlying the glueball spectrum has measurable
physical consequences across approximately 61 orders of magnitude of
spatial scale.

\begin{table}[t]
\caption{Perpendicularity test results for $N_{\text{fil}} = 3$
cosmic web nodes. All four tests significant at $p < 10^{-36}$.
Data from Tempel \textit{et al.}~\cite{Tempel2014}.}
\label{tab:perptest}
\begin{ruledtabular}
\begin{tabular}{lccc}
Test & $r$ & $p$-value & Result \\
\hline
Perpendicularity vs richness & 0.494 & $3.3\times10^{-59}$ & Signif. \\
Top vs bottom quartile       & $\Delta=0.238$ & $2.6\times10^{-40}$ & $1.7\times$ \\
$\sin(\theta)$ within $N=3$  & 0.512 & $3.9\times10^{-64}$ & Signif. \\
3D orthogonality vs richness & 0.397 & $8.0\times10^{-37}$ & Signif. \\
\end{tabular}
\end{ruledtabular}
\end{table}

%% ================================================================
\section{Relationship to Existing Work}
\label{sec:existing}

Douglas~\cite{Douglas2026} identifies probabilistic lattice gauge
theory, tensor renormalization group methods, and supersymmetric
analogs as the main active approaches to the mass gap problem. The
present work differs from all three in that it derives the glueball
spectrum from geometric mode structure rather than from dynamical
field equations. The $SU(3)$ symmetry is derived from the
perpendicular frame [$U(3)/U(1)$] rather than postulated. The mass
gap follows from colour neutrality and Pythagoras rather than from
confinement dynamics.

The frame field reformulation proposed here is related to but distinct
from the vielbein (tetrad) formalism used in general relativity. In
GR, the vielbein encodes the metric; here, the perpendicular frame
encodes the gauge structure of the strong interaction. The glueball
spectrum does not appear in standard vielbein formulations because
those formulations do not identify the mode indices
$\{1, \sqrt{2}, 2\}$ as the characteristic right-angled lengths of
the frame geometry.

%% ================================================================
\section{Open Problems}
\label{sec:open}

\textit{(i) Absolute mass scale.}---The present derivation produces
ratios only. The absolute scale $m(0^{++}) = 1.710$~GeV requires
knowledge of the QCD confinement scale $\Lambda_{\text{QCD}}$, which
we have not derived from geometry. We note that the one-loop beta
function coefficient $b_0 = 11$ is fixed by the $SU(3)$ derivation,
which reduces the problem to determining a single coupling parameter.

\textit{(ii) Quantum number assignment.}---The mapping between
perpendicular mode indices and glueball quantum numbers ($J^{PC}$) is
guided by lattice QCD ordering and symmetry considerations. A
first-principles derivation of which modes carry which quantum numbers
is needed.

\textit{(iii) Wightman axioms.}---A complete solution to the
Millennium Problem requires constructing a rigorous quantum field
theory satisfying the Wightman or Osterwalder-Schrader
axioms~\cite{Jaffe2000}. The frame field formulation is structurally
compatible with these requirements but the formal proof has not been
written.

\textit{(iv) Self-interaction corrections.}---The pure geometric
spectrum matches lattice QCD at 1.0\% mean error. The residual
discrepancies may arise from the non-Abelian self-interaction, which
redistributes energy from asymmetrically excited modes toward
equilibrium. A perturbative correction proportional to the strong
coupling constant at the glueball scale may further improve agreement.

%% ================================================================
\section{Conclusion}
\label{sec:conclusion}

The six lightest glueball states in pure $SU(3)$ Yang-Mills theory
have mass ratios given by:
\begin{equation}
\frac{m(\text{state})}{m(0^{++})} =
\frac{\sqrt{n_1^2 + n_2^2 + n_3^2}}{\sqrt{3}}
\label{eq:main}
\end{equation}
where $(n_1, n_2, n_3)$ are mode indices drawn from
$\{1, \sqrt{2}, 2\}$---the edge, face diagonal, and double edge of a
unit cube. Zero free parameters. Six states. Mean error 1.0\%. The
glueball mass spectrum is the Pythagorean theorem applied to three
perpendicular directions with right-angled mode indices.

\begin{acknowledgments}
This work builds on the Holographic Cube Lattice framework
(Ref.~\cite{Anarchos2026a}) and the empirical junction angle
correlation (Ref.~\cite{Anarchos2026b}). The perpendicularity test
uses data from the Tempel \textit{et al.} SDSS filament
catalogue~\cite{Tempel2014}. No funding was received. Analysis
pipeline available at
\url{https://github.com/VelaAnarchos/pond-model}.
\end{acknowledgments}

\bibliography{glueball_refs}

\begin{thebibliography}{12}

\bibitem{Clay2000}
A.~Jaffe and E.~Witten,
\textit{Quantum Yang-Mills Theory},
Official Millennium Problem description, Clay Mathematics Institute
(2000).

\bibitem{Douglas2026}
M.~R.~Douglas,
The Yang-Mills Millennium problem,
\href{https://doi.org/10.1038/s42254-025-00909-2}{Nat. Rev. Phys.
\textbf{8}, 86 (2026)}.

\bibitem{Morningstar1999}
C.~J.~Morningstar and M.~Peardon,
Glueball spectrum from an anisotropic lattice study,
\href{https://doi.org/10.1103/PhysRevD.60.034509}{Phys. Rev. D
\textbf{60}, 034509 (1999)}.

\bibitem{Chen2006}
Y.~Chen \textit{et al.},
Glueball spectrum and matrix elements on anisotropic lattices,
\href{https://doi.org/10.1103/PhysRevD.73.014516}{Phys. Rev. D
\textbf{73}, 014516 (2006)}.

\bibitem{Athenodorou2021}
A.~Athenodorou and M.~Teper,
The glueball spectrum of SU(3) gauge theory in 3+1 dimensions,
\href{https://doi.org/10.1007/JHEP11(2021)172}{J. High Energy Phys.
\textbf{2021}, 172 (2021)}.

\bibitem{Jaffe2000}
A.~Jaffe and E.~Witten,
\textit{Quantum Yang-Mills Theory},
in \textit{The Millennium Prize Problems}, edited by J.~Carlson,
A.~Jaffe, and A.~Wiles (American Mathematical Society, Providence,
2006).

\bibitem{Tempel2014}
E.~Tempel, R.~S.~Stoica, V.~J.~Mart\'{i}nez, L.~J.~Liivam\"{a}gi,
G.~Castellan, and E.~Saar,
Detecting filamentary pattern in the cosmic web: a catalogue of
filaments for the SDSS,
\href{https://doi.org/10.1093/mnras/stt2454}{Mon. Not. R. Astron. Soc.
\textbf{438}, 3465 (2014)}.

\bibitem{Anarchos2026a}
V.~Anarchos,
\textit{Holographic Cube Lattice Framework: Geometric Unification via
the Bidirectional Reflection Tensor, Boundary Plane Transitions, and
the Quantum Memory Matrix},
Working Paper v4 (2026).

\bibitem{Anarchos2026b}
V.~Anarchos,
\textit{Cosmic Web Junction Angle Predicts Node Richness: A Test of
the Substrate Confluence Prediction},
viXra ref 17906383 (2026).

\end{thebibliography}

\end{document}
